\chapter[Einweihungen]{Einweihungen}
\quellseite{12}
Am 24. Sonnt. n. Trin. weihte die Gemeinde zu Hantinson, N. Dak., ihr neuerbautes Gotteshaus (20×36 Fuß) dem Dienste Gottes. Prediger: PP. E. Hink und D. Glöter. H. Körber.

Am 3. Sonnt. d. Adv. weihte die Zions-Gemeinde bei Vernon, Tex., ihre neue Kirche (24×36 Fuß) dem Dienste Gottes. Prediger: P. F. M. Rudi und Unterzeichneter (englisch). Den Weihact vollzog H. F. Hellmann.

Am 3. Sonnt. d. Adv. weihte die Immanuels-Gemeinde zu Stark, Nebr., ihre neue Kirche (36×82×20, Thurm 96 Fuß) dem Dienste Gottes. Prediger: PP. Lang, Möllering (englisch) und Tiemann. Den Weihact vollzog G. Seesto.

Am 3. Sonnt. d. Adv. weihte die St. Paulus-Gemeinde zu Unionville, Mich., ihre neue Kirche (36×72×16, Thurm 74 Fuß) dem Dienste Gottes. Prediger: PP. Mayer, Born und Grüber (englisch). Das Weihegebet sprach F. Wilkening.

Am 4. Sonnt. d. Adv. weihte die St. Johannes-Gemeinde zu Lone Wolf, Okla., ihr neues Kirch- und Schulgebäude dem Dienste Gottes. Es predigte Unterzeichneter und P. Lehenbauer (englisch). B. D. Richter.

Am 9. December weihte die Zions-Gemeinde bei Ainsworth, Nebr., ihre neue Schule dem Dienste Gottes. Es predigte C. Just.

Am 3. Sonnt. d. Adv. weihte die Dreieinigkeits-Gemeinde zu Peoria, Ill., ihre Schule ein. Es predigten die Professoren: Geo. Wegner und F. Lindemann. D. L. Hohenstein.
